\color{unidarkgreen}{\subsection{Логика отключения}}
\color{black}

\begin{enumerate}[label=\arabic{section}.\arabic{subsection}.\arabic*,  align=left, labelsep=4pt, leftmargin=0pt, itemindent=57pt]

\item
Устройство формирует выходное воздействие на \textit{«Откл. осн. Защ.»} по схеме «ИЛИ» в следующих случаях:
\begin{itemize}
\item Срабатывании ДТЗ;
\item Срабатывании ПДЗР;
\item Срабатывании отключающих ступеней ГЗ ШР.
\end{itemize}

\item
Устройство формирует выходное воздействие на \textit{«Откл. рез. Защ.»} по схеме «ИЛИ» в следующих случаях:
\begin{itemize}
\item Срабатывании МТЗ;
\item Срабатывании ТЗНП;
\item Срабатывании ТЗОП.
\end{itemize}

\item
Устройство формирует выходное воздействие на \textit{«Откл. защ.»} по схеме «ИЛИ» в следующих случаях:
\begin{itemize}
\item Срабатывании ДТЗ;
\item Срабатывании ПДЗР;
\item Срабатывании отключающей ступени ГЗ;
\item Срабатывании МТЗ;
\item Срабатывании ТЗНП;
\item Срабатывании ТЗОП;
\item Срабатывании ТЗ;
\item Срабатывание отсечного клапана;
\item Срабатывание предохранительного клапана.
\end{itemize}

\item
Устройство формирует выходное воздействие на  \textit{«Откл. ШР.»} по схеме «ИЛИ» в следующих случаях:
\begin{itemize}
\item Срабатывании ДТЗ;
\item Срабатывании ПДЗР;
\item Срабатывании отключающей ступени ГЗ;
\item Срабатывании МТЗ;
\item Срабатывании ТЗНП;
\item Срабатывании ТЗОП;
\item Срабатывании ТЗ;
\item Срабатывание отсечного клапана;
\item Срабатывание предохранительного клапана;
\item Внешнее отключение.
\end{itemize}


\begin{figure}[H]
\centering
\includegraphics[width=0.7\textwidth,height=0.7\textheight,keepaspectratio]{\fbpath/9901 Логика отключения/img/lo.pdf}
\captionof{figure}{Функциональный блок логики отключения ШР}\label{pict_lo:LO}
\end{figure}

\end{enumerate}